\documentclass[11pt]{article}
\usepackage{amsmath,amssymb,amsthm,mathtools,hyperref}

\begin{document}
\section*{Integration of the mitochondrial endosymbiont}
\label{prob:306}
\label{rank:198}
\par\smallskip

\noindent\textbf{Domain.} Endosymbiosis; integration and individuality. \quad \textbf{Retained status/source context.} Open historical and comparative integration question. Endosymbiotic gene transfer is established; its existence is not itself an unresolved mechanism.

\paragraph{Definitions and assumptions.} At an ancestral stage $t$, let $Z_t=(g_t,a_t,c_t)$ describe transferred functional gene fraction $g_t\in[0,1]$, metabolic autonomy $a_t\in[0,1]$ under a specified resource test, and conflict burden $c_t\geq0$ as a normalized reproductive-output cost. A finite-state evolutionary model has transition matrix $K_\theta$ per declared evolutionary step. Integration set $B$ requires jointly specified host dependence, coordinated reproduction and constrained conflict; no single gene fraction defines individuality. Data $D$ include organelle/nuclear gene trees, targeting evidence and comparative dependence measurements with uncertainty.

\paragraph{Problem statement.} Determine whether compatible histories impose necessary integration orderings and transition constraints:
    \[\mathbb{P}_\theta(\tau_B\leq H\mid Z_0=z_0),\qquad
    \tau_B=\inf\{t:Z_t\in B\}.\]
    Bound this first-passage probability over the data-compatible parameter set and identify features shared by all sufficiently supported histories. The time-step convention and starting state $z_0$ must be stated.

\paragraph{Additional binding assumptions and resolution criterion.} For a finite-state model, specify the finite grid or enumerated allowed values of the apparently continuous descriptors and the error of that coarse representation. The transition step is an evolutionary-event convention unless calibrated to physical time. Fix the initial state and full transition law, and report the positive-probability no-integration outcome by $H$. The first-passage calculation is standard for a known matrix; historical calibration and shared necessary features are the research targets.

\paragraph{Instance definitions and transition notation.} Before fitting or testing, declare the biological lineage or population unit, environmental domain, finite state space $S$, parameter space $\Theta$, and observation and sampling law linking latent histories to $D$. Specify the reference gene universe and denominator for $g_t$, including treatment of duplication, loss, replacement and nonfunctional insertions; functional transfer must be distinguished from DNA relocation alone. Specify the environment, readout, duration and normalization of the resource test defining $a_t$, and the reproductive unit, comparison condition, reference output and positive normalization defining $c_t$. Declare the associated uncertainty models. These are requirements for an instance, not supplied measurement protocols or claims that modern comparative systems replicate ancestral states. Use the row convention $K_\theta(x,y)=\mathbb{P}_\theta(Z_{t+1}=y\mid Z_t=x)$, with nonnegative entries and $\sum_{y\in S}K_\theta(x,y)=1$. The stipulated chain is Markov and time homogeneous on $S$; finite discretization alone does not establish the biological adequacy of those assumptions. Declare the treatment of extinction or host loss and identify the approximation metric for the coarse representation, including its implications for the first-passage bounds and historical conclusions.

\paragraph{First-passage conventions.} Evolutionary-event time is $t\in\mathbb{Z}_{\geq0}$ and the declared horizon is $H\in\mathbb{Z}_{\geq0}$. Set $\inf\varnothing=\infty$; the displayed definition gives $\tau_B=0$ if $z_0\in B$. The binding positive-probability requirement restricts admitted parameter values to those satisfying $\mathbb{P}_\theta(\tau_B>H\mid Z_0=z_0)>0$; it is an assumption, not a consequence of having a finite transition matrix. Report this complement probability as well as the entry probability. In particular, $z_0\in B$ is excluded by that restriction. State whether the declared transition law makes $B$ absorbing. The displayed target concerns first entry and does not itself assert persistence after entry.

\paragraph{Integration and supported-history specifications still required.} Membership in $B\subseteq S$ must be determined by the stated coordinates and must implement the joint dependence, reproduction-coordination and conflict criteria. Whether such a biologically adequate rule exists on $(g,a,c)$ remains unspecified; additional state variables would change the stated model and are not selected here. Declare the data-compatibility rule and its uncertainty interpretation separately from the rule selecting sufficiently supported latent histories. The intended history population or law, any conditioning on integration or event occurrence, and the meaning of necessity for candidate orderings remain to be specified. Declare the candidate events and treatment of ties and absent events without replacing the requirement of features shared by all sufficiently supported histories by a weaker statement. Distinguish restrictions imposed by the model from those supported by data. The catalogue's empty-compatible-set rejection rule and informative-bound requirements continue to apply.

\paragraph{Alternative formulations.} (i) A historical reconstruction uses gene-tree reconciliation to constrain transfer chronology. (ii) A comparative model tests how modern host--symbiont dependence predicts stability without treating those systems as ancestral replicas. (iii) An infinite-population invasion model asks when selfish organelle variants increase, while finite-population conflict persistence needs absorption or reinfection assumptions. These are complementary, not equivalent.

\paragraph{Mathematical-status qualification.} The restricted calculation, identity, obstruction or construction identified in this entry is not itself a new open conjecture. The residual target is the explicitly delimited extension, identification, implementation or independent biological validation; an already known special case does not satisfy that residual target.

\paragraph{Obstructions and partial progress.} Organelle-to-nucleus transfer and its continuing genomic consequences are well documented. The supplied review frames transitions in individuality but does not determine a mitochondrial integration trajectory. Gene presence alone cannot date acquisition of functional targeting or loss of autonomy. Dependence may differ across environments, and conflict suppression can involve host or organelle changes. A reduced integration model is useful only if its coordinates are independently measurable and alternative histories can be distinguished or honestly bounded.

\paragraph{Sources (supplied provenance; not all rechecked).} \begin{itemize}
\item West et al., Major evolutionary transitions in individuality (2015). \url{https://pmc.ncbi.nlm.nih.gov/articles/PMC4547252/} (Indexed review abstract checked via PubMed 25964342; direct PMC open encountered a browser challenge.)
\item Timmis et al., Endosymbiotic gene transfer: organelle genomes forge eukaryotic chromosomes (2004). \url{https://pubmed.ncbi.nlm.nih.gov/14735123/} (Indexed review checked; established transfer context, not unique integration chronology.)
\item Indexed abstract for the supplied West et al. review. \url{https://pubmed.ncbi.nlm.nih.gov/25964342/}
\end{itemize}


\section*{Binding catalogue conventions}
This edition revises both supplied catalogues and preserves original problem numbers 1--330. It adds one globally unique label \verb|\label{rank:x}| per problem, where smaller $x$ means greater provisional difficulty. A rank label is an anchor at that problem's section: \verb|\ref{rank:x}| returns the section number, not the literal difficulty number. The visible \texttt{rank:x} text and the index give the difficulty number itself.

The entries are \emph{parameterized mathematical research schemes}, including empirical prediction, causal identification, historical reconstruction, theory and engineering. They are not 330 certified unsolved mathematical conjectures. A rigorous instance fixes the objects and restrictions below; changing those objects changes the question. Established identities, finite calculations, non-identification constructions and solved special cases are not being posed anew as open. Some residual tasks concern informative bounds, biological realization or external validation rather than a new theorem.

Every main definition and problem statement received a mathematical editorial review. Each entry now has an additional binding-conditions paragraph; direct corrections are separately recorded in the accompanying audit. The alternative formulations inherit the applicable conventions, but are different targets unless equivalence is proved. Where an old broad phrase conflicts with an explicit restriction below it, the explicit restriction controls.

\textbf{Literature boundary.} Thirteen selective primary-literature checks cover eighteen entries in this revision; their access levels and scope are recorded explicitly. The supplied bibliography and its earlier source-access descriptions are retained as provenance, not represented as newly verified. Neither this review nor a source's description proves the absence of later solutions. The date of this edition is 13 October 2026; there is no claim of exhaustive current-status certification for all 330 entries.

\subsection*{Binding conventions for a concrete instance}
\paragraph{Objects, quantifiers and non-vacuity.}
An instance declares measurable state, input, output and parameter spaces; a fixed candidate family; admissible initial laws, interventions and environments; a finite horizon or specified asymptotic regime; observation and sampling laws; and performance metrics, positive normalization scales and positive tolerances. These are chosen before fitting or testing, not adjusted after seeing failures. Shared parameters must genuinely be shared across the declared validation domain. Finite dimension alone does not restrict arbitrary real-number encodings: specify a finite library, bounded description length, or regularity/complexity bounds. Exclude constant or no-signal solutions through a fixed informative-output or baseline-improvement requirement where appropriate. A universally quantified result cannot be established by choosing one trivial favorable instance.

\paragraph{Probability laws and data.}
Every expectation or probability refers to the specified stochastic law, including initial uncertainty, process noise, observation noise and sampling as appropriate. A deterministic equation is not a probability model without a declared random input or initial law. ``Observed'' targets in population-level formulas mean population observables $g_j^*$, not noisy estimates asserted to be exact. Let $\widetilde g_j$ be independent validation estimates and suppose
\[
\mathbb{P}\{\lvert\widetilde g_j-g_j^*\rvert\le r_j\text{ for every }j\}\ge1-\alpha.
\]
Then an uncertainty-aware sufficient certification rule for positive target tolerances $\varepsilon_j$ is
\[
\max_j\frac{|\widehat g_j-\widetilde g_j|+r_j}{\varepsilon_j}\le1.
\]
On that simultaneous confidence event it implies every population error is at most its tolerance, by the triangle inequality. Vector/norm versions use corresponding valid simultaneous bounds. Training uncertainty must also be included when the claim averages over repeated training data. Evaluation on finitely many conditions does not establish a continuous-domain supremum without covering and regularity bounds.

\paragraph{Dynamics and well-posedness.}
For ODEs and SDEs give coefficient domains, initial laws and existence/uniqueness conditions in the intended solution sense, or keep the entire admitted solution set. For stochastic jump models specify nonnegative off-diagonal rates, diagonal normalization, rate orientation and non-explosion. For PDEs specify spatial domains, boundary/initial conditions, conservation or source terms, and a weak or classical solution class. Moving-domain concentration equations include dilution/transport terms. A stationary law is used only for a time-homogeneous augmented model under existence assumptions; a stationary approximation to a driven system must be justified. Absorbing extinction states are retained. Survival-conditioned, quasi-stationary and unconditional laws are distinct and must not be substituted for one another.

\paragraph{Identification and causal meaning.}
Two parameter values are observationally equivalent when they induce the same full distribution under all allowed observation/intervention designs, not merely the same likelihood value for one sample. A target is identified only if constant on the equivalence class. Gauge symmetries, label switching and reference conventions are quotiented or fixed. Finite-data compatible sets use a declared confidence, likelihood or discrepancy rule; an empty set signals model rejection. Bayesian credible sets are not automatically confidence sets. A prior or regularizer does not create data identification. Bounds intended to resolve a question must be informative at a predeclared width or decision threshold.

A causal contrast uses well-defined intervention versions and a stated causal model. Random assignment, or explicit consistency, overlap and exchangeability assumptions, is required for causal interpretation from observations. Interference and time-varying confounding are represented when present. Fixing a mediator by intervention differs from conditioning on its observed value. Marginal potential-outcome laws do not determine individual cross-world effects or an arbitrary joint coupling. Ethically admissible designs, rather than unconstrained interventions, define the experimental menu.

\paragraph{Metrics, units, optimization and limits.}
Specify the sampling or reference measure for densities and integrals, the log base, and dimensionless arguments of logarithms. Total variation is $\sup_A|P(A)-Q(A)|$; Wasserstein distances require the declared ground metric and appropriate moment conditions; KL divergence uses its stated absolute-continuity convention. All normed readouts are scaled consistently, and error probabilities use $0<\alpha<1$. A ratio has a positive denominator or an explicit undefined/zero convention. Thermodynamic entropy production is nonnegative in expectation under the relevant hypotheses; individual stochastic realizations need not be. An infimum is not an attained minimum without compactness/coercivity and continuity or another existence argument. Fix the order of large-population, long-time and vanishing-mutation limits. Stability claims declare perturbation space and symmetries.

\subsection*{How the provisional difficulty order was assigned}
Four editorial components are scored from 1 (lower burden) to 5 (higher burden): (a) formal or inferential difficulty, (b) access to evidence and identification burden, (c) mechanistic/computational integration, and (d) independent validation and transfer. Their unweighted total orders the entries. For selected entries in the highest bands, an explicit editorial priority breaks equal totals; all other ties are broken by original problem number. Equal totals are \emph{not} evidence of reliably different true difficulties. Scores are judgments, not measurements, completion-time estimates or consensus rankings.

The object ranked is an informative, nondegenerate resolution of the \emph{scoped} scientific target, with the stated budget and validation domain. A sharp useful impossibility result can qualify; an already known identity, vacuous bound, arbitrary favorable model or lookup table cannot. Rank 330 is merely the least difficult relative assessment in this catalogue, not a claim of an easy biological solution. Broad historical titles do not automatically outrank precisely formulated mechanistic tasks.



\end{document}
